\begin{statement}{D.56}{Proved}{The local section theorem for full formations containing two primes}
Let \(k/\mathbb Q_p\) be finite and let \(\mathcal C\) be a nonempty
class of finite groups closed under subgroups, quotients, and
extensions.  Suppose \(C_p,C_\ell\in\mathcal C\) for some prime
\(\ell\ne p\).  Every continuous section of the geometrically
pro-\(\mathcal C\) fundamental group of a smooth geometrically
connected hyperbolic curve over \(k\) is induced by a rational
point or belongs to a rational cusp packet.  For a proper curve
only the point alternative occurs.  In particular this holds
for the maximal geometrically pro-\(\{p,\ell\}\) quotient.
There is no bound on derived length; no one-prime assertion or
fixed finite-derived assertion is included.
\end{statement}
\begin{proof}
We first take \(X/k\) proper and hyperbolic; the affine case is
proved in the final paragraph.  We check every change to D.54,
including the case \(\ell=2\), rather than assuming that the
solvable proof is automatically compatible with an arbitrary
formation.

\paragraph{1. The exact finite-group closure argument.}

Write $\Sigma(\mathcal C)=\{q:C_q\in\mathcal C\}$.  Closure under
extensions implies closure under finite direct products.  It also
implies that every finite solvable group whose prime divisors lie
in $\Sigma(\mathcal C)$ belongs to $\mathcal C$, by a composition
series with cyclic prime-order factors.  By Burnside's theorem,
every finite $\{p,\ell\}$-group therefore lies in $\mathcal C$.
Conversely, if $G\in\mathcal C$, Cauchy's theorem and subgroup
closure imply that every prime dividing $|G|$ lies in
$\Sigma(\mathcal C)$.

If $Y\to X$ is a geometric finite Galois cover with group in
$\mathcal C$, and a finite Galois cover of $Y$ has group in
$\mathcal C$, its normal closure over $X$ again has group in
$\mathcal C$.  Over $Y$ the normal closure has group a subgroup of
a finite product of the conjugate covering groups; its quotient
is the group of $Y/X$.  The two closure properties prove the
assertion.  The same proof holds for finite function-field
extensions.  Thus, for a characteristic open subgroup $V$ of
$\Delta=\pi_1(X_{\bar k})^{\mathcal C}$ and an actual section $s$,
the group $V s(G_k)$ is exactly the geometrically pro-$\mathcal C$
fundamental group of its arithmetic cover.  There is no loss of
covers in making this identification.

Fix the single set $\Sigma_0=\{p,\ell\}$ on every cover.  All the
pro-$\Sigma_0$ covers used in D.54, Step 2, belong to $\mathcal C$.
\cite[Theorem F(i), Proposition 2.4(iv)]{MT2024}
therefore gives the unique fixed pro-$\Sigma_0$ point on every
proper hyperbolic section neighborhood.  Pullback of the target
$\Sigma_0$ tower gives the same equivariant map as in D.54;
uniqueness in the target proves compatibility.  This argument
does not require the finite quotient $\Delta/V$ to be solvable.
A compatible point of type one gives a rational point by Tate's
fixed-field theorem exactly as before.  It remains to exclude a
nonclassical fixed valuation.

\paragraph{2. All collar and patching groups remain in the formation.}

The collar construction of D.54, Step 3, uses cyclic groups of
order $\ell$, normal closures of towers of such groups, and
geometrically pro-$\Sigma_0$ stable-model domination.  Each belongs
to $\mathcal C$ by Paragraph 1.  This also covers the separating-edge
construction: its repeated cyclic covers and their normal closure
are $\ell$-groups.  In the positive-genus component case an
unramified cyclic cover of order $\ell$ exists because the residue
characteristic is $p\ne\ell$.  The orders at the two marked points
are congruent to $+1$ and $-1$ modulo $\ell$; this gives the
opposite inertia characters required by the invariant annular
gluing, also when $\ell=2$.

In the equivariant patching step let the specified exterior
Galois group be $G\in\mathcal C$.  The construction produces an
actual $G$-torsor on the auxiliary curve; each connected component
has group a subgroup of $G$.  Its normal closure over the base
curve has group in $\mathcal C$ by Paragraph 1.  Thus replacing
``solvable'' with ``in $\mathcal C$'' preserves the entire patching
argument, including the group action and the algebraization.

For a general finite Galois extension $F/\bar k(X)$ with group
$G\in\mathcal C$, every horizontal ramification index divides
$|G|$.  Choose each auxiliary root order $m$ to be the least common
multiple of the relevant indices.  Every prime divisor of $m$
lies in $\Sigma(\mathcal C)$.  The auxiliary cover
\[
 Z:\quad y^\ell=T(T-1)/(T-\alpha),\qquad
 Z'=Z\times_{J,[m]}J
\]
has cyclic initial group $C_\ell$ and finite abelian division
group $J[m]$.  Over $\bar k$ the latter is a product of cyclic
groups with prime divisors among those of $m$, and hence belongs
to $\mathcal C$.  Its normal closure, all conjugates, all finite
composita, and the patched covers stay in $\mathcal C$.  No prime
outside $\Sigma(\mathcal C)$ is introduced by the division step.
Horizontal indices still divide $\ell$, for the characteristic-zero
local-field reason in D.54.

\paragraph{3. The case $\ell=2$ in the auxiliary curve.}

If $\ell$ is odd, the auxiliary computation is literally D.54,
Step 4.  If $\ell=2$, then $p$ is odd.  The curve
\[
                    y^2=T(T-1)/(T-\alpha)
\]
is a smooth genus-one double cover, with four distinct branch
points $0,\alpha,1,\infty$.  For $a=v_p(\alpha)>0$, the same
semistable model has two rational components and two parallel
annular edges of length $a$.  One component is obtained from the
outer coordinate $T$, the other from $T=\alpha u$; their reduced
equations are $y^2=T-1$ and $y^2=-u/(u-1)$.  At the common end
there are two distinct points, because $p\ne2$.  The skeleton is
a circle of total length $2a$.

This model is semistable rather than stable as an unmarked
genus-one curve.  Stability is not used in the auxiliary division
construction.  Its Jacobian is the same totally degenerate
elliptic curve and has a toric uniformization
\cite[Sections 4.1--4.2 and 4.5.2]{DE_BakerRabinoff2015}
$\mathbb G_m^{\mathrm{an}}/\Lambda$.  The cycle $e_1-e_2$ is a
primitive torus character; its Abel--Jacobi pullback on either
edge has slope one, with the choice of edge orientation.  The
same calculation as in D.54 gives
\[
 u_j=c_jT(1+h_j),\qquad
 |h_j|\le p^{-\min\{v_p(T),a-v_p(T)\}}.
\]
After division by $[m]$, the period adjustment on the clopen
$\Lambda/\Lambda^m$ pieces and the torsion-freeness of $\Lambda$
produce an analytic $m$-th root of $u_j$ on every entire inverse
edge.  On the middle collar with
$D>v_p(m)+1/(p-1)$ and $a>2D$, the log--exp root of $1+h_j$
then gives an $m$-th root of $T$.  The finite groups here are
still a cyclic group of order two followed by a subgroup of
$J[m]$.  Thus the case $\Sigma_0=\{p,2\}$ needs no third prime.

\paragraph{4. The literal completion inclusion.}

Let $B/\bar k(X)$ now be the maximal pro-$\mathcal C$ birational
extension, let $L$ be the universal geometrically pro-$\mathcal C$
\'etale extension, and let $H=\widehat L$ for the proposed
nonclassical valuation.  The proof of D.54, Step 5, applies
verbatim with the group checks of Paragraphs 1--3.  First every
finite $\mathcal C$-Galois extension with all horizontal indices
dividing $\ell$ embeds literally in $H$, using the collar lemma,
equivariant patching, and the complete-field-valued point
transfer.  General finite $\mathcal C$-Galois extensions are then
treated by the auxiliary $m$-division covers.  Their normal
closures have indices dividing $\ell$ and group in $\mathcal C$,
so the first step supplies their function fields in $H$.
Normality of each prescribed $F/\bar k(X)$ makes its resulting
embedding have the literal chosen image $F\subset H$.
Consequently $B\subset H$.

Exactly as in D.54, the map of decomposition groups
$D_w(B/k(X))\to D_w(L/k(X))$ is an isomorphism: prolongation
transitivity gives surjectivity, while an element in the kernel
fixes the completion $H$ pointwise and therefore fixes $B$.
The local section lifts to a geometrically pro-$\mathcal C$
birational section.  \cite[Proposition 1.7]{DE_Hoshi2012} applies because
$p\in\Sigma(\mathcal C)$, and places its whole image in the
horizontal decomposition group of a rational point.

\paragraph{5. The final contradiction uses any prime of the formation.}

This is the only alteration to D.54, Step 6, besides the group
names.  Choose any prime $q\in\Sigma(\mathcal C)$, for example
$q=p$.  The field $B$ is $q$-closed: if $b\in B$, put $b$ in a
finite $\mathcal C$-Galois extension of $\bar k(X)$ and adjoin
$q$-th roots of all conjugates of $b$.  Since $\bar k$ contains
$\mu_q$, the resulting normal extension has an elementary
abelian $q$-group as its relative Galois group; its full group
belongs to $\mathcal C$.  Thus all $q$-th roots of $b$ already
lie in $B$.

Put $M=B^{s(G_k)}$.  Both the vertical valuation $w_1$ and the
horizontal valuation $w_2$ at the rational point have unique
prolongation to $B/M$, because their decomposition group is all
of $\operatorname{Gal}(B/M)=s(G_k)$.  Hence
\[
                          B\cap M^{h_i}=M\quad(i=1,2).
\]
Indeed a henselization is the fixed field of the decomposition
group in a separable closure; restriction of that decomposition
group onto $B$ is all of $\operatorname{Gal}(B/M)$.
If $u\in M$ and
\[
                       w_i(u-1)>2w_i(q),
\]
Hensel's lemma for $Z^q-u$ at $Z=1$ gives a root in $M^{h_i}$.
Every root lies in $B$ by $q$-closedness, so the root lies in $M$.
This argument requires no assumption that $\mu_q\subset k$ or
$\mu_q\subset M$.

Choose a uniformizer $d$ of $k$.  It is not a $q$-th power in $k$,
and cannot become one in $M$ because $M\cap\bar k=k$.  Let $t$ be
a horizontal uniformizer at the rational point.  Choose $N$ with
$Nw_1(p)>w_1(t)$ and set $z=p^N/t$.  Then $w_1(z)>0$ and
$w_2(z)<0$.  For large positive integers $r$, set
\[
                  u_r=d+(1-d)\frac{z^r}{1+z^r}.
\]
We have $w_1(u_r/d-1)\to+\infty$ and
$w_2(u_r-1)\to+\infty$.  Applying the preceding root criterion
to $u_r/d$ for $w_1$ and to $u_r$ for $w_2$ makes both elements
$q$-th powers in $M$.  Their quotient is $d$, a contradiction.
The nonclassical fixed valuation is therefore impossible.

\paragraph{6. The affine passage and scope.}

The affine compactification argument of D.54, Step 7, only needs
cofinal characteristic finite covers with all peripheral orders
arbitrarily large.  The pro-$p$ covers used there belong to
$\mathcal C$ because $C_p\in\mathcal C$.  Riemann--Hurwitz gives
hyperbolic compactifications; the proper result, common refinements,
and compactness give the compatible point-or-cusp decomposition
in the original affine quotient.  This completes the proof.

The result uses a full formation, hence unrestricted derived
length.  It neither contradicts the nongeometric one-prime
sections constructed in D.12/D.52 nor settles a fixed finite
derived stage.  Its most restrictive specialization is the
geometrically pro-$\{p,\ell\}$ quotient.  It makes no assertion
about a pro-$\ell$ quotient when $p\notin\Sigma$, or about a
geometrically pro-$p$ quotient with no second prime.
\end{proof}
